% 2.1.5
% Ficheiro incluído por main.tex

\subsection[Samba: acesso à partilha SRII com Kerberos]{Samba: acesso à partilha SRII com Kerberos}

\passo{Passo 1 - Limpar a cache do Kerberos}

\begin{lstlisting}[style=bashstyle]
kdestroy
klist
\end{lstlisting}

\passo{Passo 2 - Listar as pastas partilhadas do servidor}

\begin{lstlisting}[style=bashstyle]
kinit aluno.teste@ISEC.LOCAL
smbclient -L //SR2-SRV.isec.local --use-kerberos=required
\end{lstlisting}

\emph{Verificação esperada:} aparece a partilha \textbf{SRII}.

\passo{Passo 3 - Preparar um ficheiro de testes}

\begin{lstlisting}[style=bashstyle]
echo "Escrito a partir do Ubuntu" > /tmp/teste-escrita.txt
chmod 644 /tmp/teste-escrita.txt
kdestroy
klist
\end{lstlisting}

\passo{Passo 4 - Testar o aluno}

\begin{lstlisting}[style=bashstyle]
kinit aluno.teste@ISEC.LOCAL
smbclient //SR2-SRV.isec.local/SRII --use-kerberos=required -c 'ls'
smbclient //SR2-SRV.isec.local/SRII --use-kerberos=required -c 'put /tmp/teste-escrita.txt aluno-escrita.txt'
kdestroy
\end{lstlisting}

\emph{Verificação esperada:} \texttt{ls} funciona; \texttt{put} falha com \texttt{NT\_STATUS\_ACCESS\_DENIED}.

\passo{Passo 5 - Testar o professor}

\begin{lstlisting}[style=bashstyle]
kinit professor.teste@ISEC.LOCAL
smbclient //SR2-SRV.isec.local/SRII --use-kerberos=required -c 'ls'
smbclient //SR2-SRV.isec.local/SRII --use-kerberos=required -c 'put /tmp/teste-escrita.txt professor-escrita.txt'
smbclient //SR2-SRV.isec.local/SRII --use-kerberos=required -c 'ls'
kdestroy
klist
\end{lstlisting}

\emph{Verificação esperada:} o professor lista e escreve; \texttt{professor-escrita.txt} aparece no segundo \texttt{ls}.

\passo{Passo 6 - Analisar resultados e discutir}

\begin{center}\small
\begin{tabularx}{\linewidth}{|>{\raggedright\arraybackslash}X|>{\raggedright\arraybackslash}X|>{\raggedright\arraybackslash}X|}
\hline
\rowcolor{azul!10}\textbf{Utilizador} & \textbf{Listar (\texttt{ls})} & \textbf{Escrever (\texttt{put})} \\
\hline
\texttt{aluno.teste} & Permitido & \texttt{NT\_STATUS\_ACCESS\_DENIED} \\
\hline
\texttt{professor.teste} & Permitido & Permitido \\
\hline
\end{tabularx}
\end{center}

Confirmar também no Windows Server:

\begin{lstlisting}[style=psstyle]
Get-ChildItem "$env:USERPROFILE\Desktop\SRII"
\end{lstlisting}

\emph{Verificação esperada:} existe \texttt{professor-escrita.txt} e \textbf{não} existe \texttt{aluno-escrita.txt}.
